\paragraph{What the enriched-space comparison does and does not identify.}
The Legendre35 and Legendre36 controls also contain the same six analytic
layers in Eq.~\eqref{eq:layer-features}. The present experiment therefore does
not show that an unenriched Legendre space succeeds where an unenriched
trigonometric space fails. Both families use the same prescribed radial
quadrature at each resolution; a Legendre trial basis does not automatically
replace it with a family-specific Gauss--Lobatto grid. Endpoint clustering of
polynomial interpolation nodes is consequently not a measured explanatory
variable in this comparison. An unenriched, rank-matched polynomial control
would require a separate experiment.

The layer functions decay as $e^{-\gamma(L-s)}$ and are concentrated near the
\emph{outer} boundary $s=L$, not the cylinder wall $s=0$. Moreover, the
homogeneous far-field operator uses $\chi=16z^2(1-z)$, with the correction in
Eq.~\eqref{eq:boundary-transform}. Its no-slip conditions do not force the
inner-wall curvature to vanish: with primes denoting $z$ derivatives,
$(\mathcal Bq)''(0)=32q(0)-[(\chi q)'']_{z=1}$, which is generally nonzero.
Thus removing these layers cannot by itself establish a chain from lost
inner-wall curvature to a changed separation angle. The large final errors
demonstrate sensitivity to the trial space and retained dimension together.
Likewise, a midpoint design matching or slightly improving the worst error
of five sampled seeds is not a theorem that randomization never helps.
